\documentclass[10pt,a4paper]{amsart}
\usepackage{amsmath,amssymb,amsthm,graphicx}
% Minimal class-compatibility preamble for a reader excerpt of
%   Bhimarthi Ravinder, "Demazure Modules, Chari-Venkatesh Modules and Fusion Products",
%   SIGMA 10 (2014), 110, DOI 10.3842/SIGMA.2014.110.
%
% The author's pinned TeX file declares no theorem environments of its own: they are
% supplied by the SIGMA publisher class sigma.cls, which is not installed here.  The
% block below is that class's theorem-environment block (same environment names and
% styles), copied mechanically so that the retained author text compiles unchanged.
% No mathematical macro or statement is introduced by this file.
\newtheorem{theorem}{Theorem}
\newtheorem*{theorem*}{Theorem}
\newtheorem{lemma}{Lemma}
\newtheorem*{lemma*}{Lemma}
\newtheorem{corollary}{Corollary}
\newtheorem*{corollary*}{Corollary}
\newtheorem{proposition}{Proposition}
\newtheorem*{proposition*}{Proposition}
\newtheorem{conjecture}{Conjecture}
\newtheorem*{conjecture*}{Conjecture}
{\theoremstyle{definition} \newtheorem{definition}{Definition}
\newtheorem*{definition*}{Definition}
\newtheorem{example}{Example}
\newtheorem*{example*}{Example}
\newtheorem{remark}{Remark}
\newtheorem*{remark*}{Remark}
\newtheorem{note}{Note}
\newtheorem*{note*}{Note}
}
\usepackage{cleveref}
\crefname{originalreference}{original reference}{original references}
\Crefname{originalreference}{Original reference}{Original references}
\makeatletter
\newlabel{f2}{{f2}{source}{}{sourcepointer-1}{}}
\newlabel{f2@cref}{{[originalreference][0][]f2}{[1][0][]source}{}{}{}}
\makeatother
\title{A surjection from the module to a fusion product and the resulting lower bound for its dimension}
\author{Bhimarthi Ravinder}
\date{doi:10.3842/SIGMA.2014.110; CC BY 4.0}
\begin{document}
\maketitle
\noindent\textit{Editorial note.} All delivered mathematical text is the author's own, copied from the pinned article source sources/sigma-2014-110-demazure-fusion/pinned-original-source.tex; no statement or proof step was rewritten, reordered or removed.\par\medskip
\noindent\textit{Editorial note.} Preamble provenance: the author's tex declares no theorem environments of its own (they come from the SIGMA publisher class); the staged sources/sigma-2014-110-demazure-fusion/retained-preamble.tex reproduces that class's theorem-environment block verbatim and is used as the preamble, with the wrapper class and base packages supplied by the assembler.\par\medskip
\noindent\textit{Editorial note.} Bibliography provenance: neither the counted statement nor its proof contains a citation (this lane rescans the whole span before writing the card), so the card carries no bibliography and no entry was reconstructed.\par\medskip
\noindent\textit{Editorial note.} Selection provenance: this card was opened under the file-based control contract (CONTROL.md), with the mechanical pre-card gate run on its statement and proof interval before the card was written.\par\medskip
\noindent\textit{Allowed standard prerequisites (editorial).}\begin{itemize}
\item The fusion product of graded modules with its notation, and the modules D(1, theta) and the evaluation module of the fundamental representation appearing as its factors. The author's sections 3 and 4.1, where the modules and the fusion product are set up before the counted statement. Fixes the target of the asserted surjection and the two factors whose dimensions occur in the inequality.
\item The defining relations of the cyclic module on the source side, and the annihilation of the generators of the two factors by twisted positive and negative root vectors. Used inside the counted proof (source lines 670-697). Supplies the relations which, by the lemma labelled f2, produce the annihilation needed to verify the presentation.
\item The statement that annihilation in the factors passes to the fusion product. Invoked inside the counted proof at source line 686 as a pointer to the lemma labelled f2 (a navigation-pointer stub). Transfers the relation satisfied by the image of the generator to the fusion product.
\end{itemize}\medskip
\ifdefined\hypertarget\hypertarget{sourcepointer-1}{}\fi
\noindent\textit{Source reference.} Original reference f2: sources/sigma-2014-110-demazure-fusion/pinned-original-source.tex, source lines 686--686. Navigation pointer only: the counted text refers to the item f2 of this same article, which is not part of this card. The pointer carries the item's own author label as its number; its content is not reproduced here.\par\medskip
\begin{proposition}
The map $\psi :\mathbf{V}_{i,k+1}\rightarrow D(1,\theta)^{* (k+1-i)} * \textup{ev}_0 V(\theta)^{*i}$ such that
$\psi(\overline{v}_{i,k+1}) = \overline{w}_{\theta}^{* (k+1-i)} * v_{\theta}^{*i}$ is well-defined and surjective
morphism of $\mathfrak{g}[t]$-modules.
In particular,
\begin{gather}
\label{dimg}
\dim \mathbf{V}_{i,k+1}\geq (\dim D(1, \theta))^{k+1-i} (\dim V(\theta))^{i}.
\end{gather}
\end{proposition}

\begin{proof}
We only need to show that $\psi(\overline{v}_{i,k+1})$ satisf\/ies the def\/ining relations of $\mathbf{V}_{i,k+1}$. But
they follow easily from the following relations:
\begin{gather*}
((h \otimes 1)-\langle (k+1)\theta, h\rangle)\big(\overline{w}_{\theta}^{\otimes (k+1-i)} \otimes v_{\theta}^{\otimes i}\big)=0,
\qquad
\forall\, h\in\mathfrak{h},
\\
(x^{-}_{\alpha}\otimes 1)^{\langle (k+1)\theta,
\qquad
\alpha^{\vee} \rangle + 1}\big(\overline{w}_{\theta}^{\otimes (k+1-i)}\otimes v_{\theta}^{\otimes i}\big)=0,
\qquad
\forall\, \alpha\in R^{+}
\end{gather*}
(which holds in $D(1,\theta)^{\otimes (k+1-i)} \otimes \textup{ev}_0 V(\theta)^{\otimes i}$) and $(h \otimes
t^s)\psi(\overline{v}_{i,k+1})=0, \forall\, s\geq 1$, $h\in\mathfrak{h}$
(which holds in $D(1,\theta)^{* (k+1-i)}*\textup{ev}_0 V(\theta)^{*i}$).
Further from Lemma~\ref{f2}, by using the relations
\begin{gather*}
(x^{+}_{\alpha}\otimes t^s)\overline{w}_{\theta}=0=(x^{+}_{\alpha}\otimes t^s)v_{\theta},
\qquad
\forall\, s\geq 0, \quad \alpha\in R^{+},
\\
(x^{-}_{\alpha}\otimes t)\overline{w}_{\theta}=\big(x^{-}_{\theta}\otimes
t^{2}\big)\overline{w}_{\theta}=0=(x^{-}_{\theta}\otimes t)v_{\theta}=(x^{-}_{\alpha}\otimes t)v_{\theta},
\qquad
\forall\, \alpha \in R^{+}\setminus\{\theta\},
\end{gather*}
which holds in $D(1,\theta)$ and $\textup{ev}_0 V(\theta)$.
\end{proof}

\end{document}
